% !TEX program = xelatex
% !TEX encoding = UTF-8 Unicode
%
% <과목 이름> --- 강의노트 국문판.  하나의 연속 문서이고 주제순이다.
% 영문판과 함께 쓰는 설정은 전부 coursenotes.sty 와 figures.tex 에 있다.
% 빌드: ./build.sh   (latexmk -xelatex, 목차 때문에 세 번)
%
% 강의를 더할 때의 규칙 (notes-en.tex 머리의 목록과 같다):
%   1. 새 내용은 주제가 맞는 대단원에 넣는다.  정말로 새 주제일 때만
%      \partdivider 로 새 대단원을 연다.
%   2. 절 번호는 대단원과 무관하게 이어진다.  다시 1부터 세지 않고, 이미 게시한
%      절의 번호는 바꾸지 않는다 (사이트 뷰어가 그 번호를 쓴다).
%   3. 그 강의의 새 내용이 시작되는 자리에 \lecturebreak[YYYY-MM-DD]{제N강},
%      바로 밑에 \sourcenote{...}, 그리고 앞부분 회차표에 한 줄을 더한다.
%   4. 절·식·그림·정리류 환경과 모든 \label 은 두 판에서 순서와 이름이 같다.
%   5. \lecdate 와 "강의의 흐름" 목록을 고친다.
%   6. \eqref/\ref 뒤의 조사는 번호를 읽는 소리를 따른다
%      (tools/ko-particles-check.py 로 확인).
%
\documentclass[11pt]{article}
\usepackage[korean]{coursenotes}

%% 그림 속 낱말 (영문판에도 같은 매크로가 있다)
\newcommand{\lblPosition}{위치 $x$}
\newcommand{\lblConcentration}{$c(x,t)$}
\newcommand{\lblSpreads}{폭은 $\sqrt{Dt}$에\\비례해 커진다}
%% figures.tex --- every diagram of the notes, shared by both editions.
%%
%% LANGUAGE-NEUTRAL ON PURPOSE.  No word is typed inside a figure: every label
%% goes through a \lbl... macro that each main file defines in its preamble
%% (English in notes-en.tex, Korean in notes-ko.tex).  The same file therefore
%% draws the same picture with the same number in both editions.  Never fork a
%% figure per language; add the new word to BOTH preambles instead.  Only the
%% \caption text lives in the language files.
%%
%% Black and grey only; curves are told apart by line style (solid, dashed,
%% dotted), never by colour.  Plotted curves are computed from the formula the
%% caption states, not sketched by hand.
%%
%% One macro per figure, \fig<Name>, in the order the figures first appear.
%% Large generated data (simulation snapshots, many-point curves) goes in its
%% own macro written by a script; say which script in a comment above it.

%% ===== (1) a diffusing profile spreads like a Gaussian =====================
%% c(x,t) = exp(-x^2/4Dt) / sqrt(4 pi D t), drawn for Dt = 0.25, 1, 4.
\newcommand{\figDiffusion}{%
\begin{tikzpicture}[x=0.9cm,y=4.2cm,>=Stealth]
  \draw[->] (-4.3,0) -- (4.5,0) node[right] {\lblPosition};
  \draw[->] (0,0) -- (0,1.32) node[left] {\lblConcentration};
  \foreach \x in {-4,-2,2,4} \draw (\x,0) -- (\x,-0.02) node[below] {\footnotesize $\x$};
  \draw[thick,samples=120,domain=-4.2:4.2]
    plot (\x,{exp(-\x*\x/1.0)/sqrt(3.14159265*1.0)});
  \draw[thick,dashed,samples=120,domain=-4.2:4.2]
    plot (\x,{exp(-\x*\x/4.0)/sqrt(3.14159265*4.0)});
  \draw[thick,dotted,samples=120,domain=-4.2:4.2]
    plot (\x,{exp(-\x*\x/16.0)/sqrt(3.14159265*16.0)});
  \node[anchor=west,font=\footnotesize] at (0.55,1.05) {$Dt=0.25$};
  \node[anchor=west,font=\footnotesize] at (1.25,0.36) {$Dt=1$};
  \node[anchor=west,font=\footnotesize] at (2.6,0.17) {$Dt=4$};
  \node[font=\footnotesize,align=center] at (-3.0,0.95) {\lblSpreads};
\end{tikzpicture}}

\endinput


\coursename{예시 과목}
\term{2027년 2학기}
\lectitle{강의노트, 제1--2강}
\lecdate{2027년 9월 6일 기준}
\notesby{수강생}

\begin{document}
\makelecheader

% ===================================================================== 앞부분
\section*{수업 정보}
\begin{itemize}[leftmargin=1.2em,itemsep=1pt,topsep=3pt]
  \item 월·수 수업. 중간·기말고사 날짜는 강의계획서에 공지된 대로.
  \item 성적 산출과 교재는 강의계획서에 공지된 대로.
\end{itemize}

\section*{강의의 흐름}
\begin{itemize}[leftmargin=1.2em,itemsep=1pt,topsep=3pt]
  \item 제1부 --- 무작위 걷기와 확산 \hfill {\small[진행 중]}
  \item 제2부 --- 정상 상태 \hfill {\small[예정: 제2강에서 예고]}
\end{itemize}

\section*{이 노트의 구성}
노트는 날짜순이 아니라 주제순이다. 오른쪽 끝에 강의 번호가 붙은 가는 선은 그
강의가 시작한 자리이고, 그 밑의 회색 줄은 그 강의의 정리가 무엇에 기대는지
적는다. 정리하면서 덧붙인 내용(빠진 단계, 검산)은 여백에 $\blacktriangleright$
표시를 단 작은 글씨로 넣었다.

\section*{회차표}
{\small
\begin{longtable}{@{}l l l r@{}}
\toprule
강의 & 날짜 & 절 & 쪽\\
\midrule
\endhead
제1강 & 2027.\,09.\,01 (수) & \S\ref{sec:walk} & \pageref{lec:2027-09-01}\\
제2강 & 2027.\,09.\,06 (월) & \S\ref{sec:diffusion} & \pageref{lec:2027-09-06}\\
\bottomrule
\end{longtable}}

{\small\tableofcontents}

% ===================================================================== 제1부
\partdivider{제1부 --- 무작위 걷기와 확산}

\lecturebreak[2027-09-01]{제1강}
\sourcenote{유인물 1과 녹음}
\section{무작위 걷기}\label{sec:walk}

\begin{definition}\label{def:walk}
직선 위의 \keyword{무작위 걷기}는 위치의 열 $x_n = \sum_{k=1}^{n} s_k$ 이고,
각 걸음 $s_k=\pm a$ 는 서로 독립이며 두 값이 나올 확률이 같다.
\end{definition}

걸음 하나의 평균은 $0$, 분산은 $a^2$ 이고 독립인 걸음의 분산은 더해지므로
\begin{equation}\label{eq:var}
  \avg{x_n} = 0, \qquad \avg{x_n^2} = n a^2 .
\end{equation}

\scribe{둘째 결과에는 걸음끼리 상관이 없다는 것만 필요하다.
$j\neq k$ 이면 $\avg{s_j s_k}=0$ 이어서 $\avg{x_n^2}$ 의 교차항이 모두 사라진다.}

% ---------------------------------------------------------------------------
\lecturebreak[2027-09-06]{제2강}
\sourcenote{유인물 2}
\section{확산 방정식}\label{sec:diffusion}

걸음 사이 시간을 $\tau$ 로 두고 $t=n\tau$, $D=a^2/2\tau$ 라 쓰면 식~\eqref{eq:var}은
$\avg{x^2}=2Dt$ 가 되고, 연속 극한에서 걷는 입자의 밀도는
\begin{equation}\label{eq:diffeq}
  \pd{c}{t} = D\,\frac{\partial^2 c}{\partial x^2}
\end{equation}
를 따른다. 모든 입자가 원점에서 출발할 때 식~\eqref{eq:diffeq}의 해는
\begin{equation}\label{eq:gauss}
  c(x,t) = \frac{1}{\sqrt{4\pi D t}}\, e^{-x^2/4Dt}
\end{equation}
이고, 그림~\ref{fig:diffusion}에 그렸다.

\begin{figure}[ht]
  \centering
  \figDiffusion
  \caption{$Dt=0.25$, $1$, $4$ 일 때의 분포 \eqref{eq:gauss} (실선, 파선, 점선).
  넓이는 $1$ 로 그대로이고 폭은 $\sqrt{2Dt}$ 로 커진다.}
  \label{fig:diffusion}
\end{figure}

\begin{caution}[두 가지 약속]\label{caut:dconv}
$D=a^2/\tau$ 로 정의해 $\avg{x^2}=Dt$ 로 쓰는 책도 있다. 식~\eqref{eq:diffeq}와
섞어 쓰기 전에 어느 약속인지 확인한다.
\end{caution}

% ===================================================================== 제2부
\partdivider{제2부 --- 정상 상태}

\lecturebreak[2027-09-08]{예정}
\section{정상 상태와 경계 조건}\label{sec:steady}

게시된 개요로 수업 전에 미리 쓴 절이다. 수업이 끝나면 칠판과 대조한다.

\phantomsection\label{lec:end}
\end{document}
